\begin{textsample}{NEG Dim 1 – 2014\_c6 – Score -10.00 – 2014\_c6/enade\_2014\_eng\_amb\_c63c648bd8.txt}  \label{ex:f1_neg_006}
CAPÍTULO 3 ANÁLISE TÉCNICA DA PROVA Este capítulo tem por objetivo apresentar o desempenho dos estudantes concluintes de Engenharia Ambiental no ENADE/2014.
Para isso, foram calculadas as estatísticas básicas da prova em seu todo ( seção 3.1.1 ), bem como as estatísticas relacionadas aos Componentes de Formação Geral ( seção 3.1.2 ) e de Conhecimento Específico da Área ( seção 3.1.3 ).
Dadas as suas características, foram analisadas, em separado, as questões objetivas ( seção 3.2 ) e as questões discursivas ( seção 3.3 ).
Para as questões objetivas, foram disponibilizados os índices de facilidade e de discriminação bisserial, também em separado, para os Componentes de Formação Geral ( seção 3.2.1 ) e de Conhecimento Específico ( seção 3.2.2 ).
De \textbf{cada} componente, uma das questões foi escolhida para exemplificar a análise gráfica, relacionando as alternativas escolhidas pelos estudantes ( inclusive o gabarito ) com o número de acertos no componente.
O Anexo I apresenta a íntegra da análise gráfica para todas as questões objetivas.
Para \textbf{cada} uma das questões discursivas, os conteúdos dos tipos mais comuns de respostas dos estudantes são apresentados e comparados com o padrão de respostas esperado ( ver Anexo VIII com o padrão de respostas ).
Tomando como base as duas questões discursivas do Componente de Formação Geral, a seção 3.3.1.6 apresenta comentários sobre a correção das respostas com respeito à Língua Portuguesa.
Na Tabela 3.1 são apresentados o tamanho da população inscrita e de presentes, e as seguintes estatísticas das notas¹⁸ : média do desempenho na prova, erro padrão da média, desvio padrão, nota mínima, mediana e nota máxima.
Tais estatísticas contemplam o total de estudantes concluintes da Área de Engenharia Ambiental inscritos e presentes à prova do ENADE/2014, tendo em vista as agregações por Grandes Regiões e o país como um todo.
As estatísticas das demais tabelas ímpares serão baseadas na mesma população da Tabela 3.1 e, portanto, do mesmo tamanho.
As três primeiras tabelas pares são desagregadas por Categoria Administrativa e Organização Acadêmica da IES.
As estatísticas dessas tabelas pares serão baseadas na mesma população da Tabela 3.2.
Em relação aos gráficos de distribuição de notas, o intervalo considerado foi de 10 unidades, aberto à esquerda e fechado à direita, com exceção do primeiro intervalo, [ 0 ; 10 ], fechado em ambos os extremos.
Para os gráficos de distribuição das ¹⁸ Uma definição dessas estatísticas pode ser encontrada no Glossário. notas das questões discursivas, foram consideradas mais duas categorias : questão em branco19 e nota zero. 3.1 ESTATÍSTICAS BÁSICAS DA PROVA Esta seção apresenta estatísticas selecionadas e histogramas da nota geral ( 3.1.1 ) e de \textbf{cada} componente : Formação Geral ( 3.1.2 ) e Conhecimento Específico ( 3.1.3 ).
São também apresentadas estatísticas selecionadas de subpopulações caracterizadas por Grande Região, Categoria Administrativa e Organização Acadêmica. 3.1.1 Estatísticas Básicas Gerais A Tabela 3.1 apresenta as estatísticas básicas da prova, por Grande Região, dos estudantes concluintes de Engenharia Ambiental.
A população total de inscritos foi de 7.329.
Destes, 6.551 \textbf{estiveram} presentes, sendo 10,6 \% o índice de não comparecimento.
A Região de maior abstenção foi a Norte ( 13,5 \% ) e a de menor abstenção foi a Sul ( 7,8 \% ).
A média das notas da prova como um todo ( nas seções seguintes serão analisados os Componentes de Formação Geral e de Conhecimento Específico ) foi 47,1, sendo que os alunos da região Centro-Oeste obtiveram a média mais baixa ( 43,9 ), e os da região Sul obtiveram a média mais alta ( 52,4 ).
As demais médias foram : 44,6 na região Norte ; 46,0 na região Nordeste ; e 46,6 na região Sudeste.
O desvio padrão para o Brasil como um todo foi 14,1, sendo o maior desvio padrão encontrado nas regiões Sudeste e Centro-Oeste ( 14,2 em ambas ) e o menor, na região Sul ( 13,1 ), indicando uma dispersão um pouco menor das notas desta última região.
A região que obteve a maior nota máxima foi a Sudeste ( 94,7 ), ao passo que a região que atingiu a menor nota máxima foi a Norte ( 79,0 ).
A mediana do Brasil como um todo foi 47,2, sendo a maior mediana obtida na região Sul ( 53,0 ), e a menor obtida na região Centro-Oeste ( 44,1 ).
A nota mínima foi zero apenas na região Sudeste, nas demais regiões foi : 7,8 na região Norte ; 13,1 nas regiões Nordeste e Sul ; e 3,8 na região Centro-Oeste.
Considerando -se as notas segundo Grande Região, observa -se que existe diferença estatisticamente significativa ao nível de 95 \% entre a menor média, atingida na região Centro-Oeste ( 43,9 ), e as médias de quase todas as regiões, menos a região Norte.
O comportamento das notas dos estudantes de todo o Brasil pode ser observado no Gráfico 3.1 que apresenta um histograma com a distribuição das mesmas.
Essa é uma distribuição unimodal com a moda no intervalo ( 40 ; 50 ].
A Tabela 3.2 apresenta as estatísticas básicas da prova, desagregadas por Categoria Administrativa e por Organização Acadêmica.
Da população total de inscritos, 5.150 são de IES Privadas e 2.179 de IES Públicas.
Em relação à Organização Acadêmica, a maior participação foi obtida por estudantes de Universidades ( 4.100 ), seguido de Faculdades ( 1.682 ) e Centros Universitários ( 1.547 ).
O nível de abstenção nas IES Públicas ( 9,5 \% ) foi menor que nas IES Privadas ( 11,1 \% ).
Entre as Organizações Acadêmicas, as Faculdades apresentaram o nível mais elevado de absenteísmo ( 11,9 \% ), acima da média nacional de 10,6 \%.
A média das notas da prova como um todo foi 47,1.
Em relação à Categoria Administrativa, os estudantes das IES Públicas obtiveram média mais alta ( 54,3 ), e os das IES Privadas obtiveram média mais baixa ( 44,0 ), que a média nacional.
Observa -se que existe diferença estatisticamente significativa entre as médias das notas das IES Públicas e Privadas.
A diferença entre as médias das regiões Sul e Centro-Oeste ( 8,5 ), a maior e a menor média, é inferior à diferença entre IES Públicas e Privadas ( 10,3 ), caracterizando uma maior diversidade administrativa do que regional.
No tocante à Organização Acadêmica, apenas as Universidades obtiveram média mais alta ( 49,8 ) que a nacional.
As demais médias foram : 44,2 nos Centros Universitários e 43,2 nas Faculdades.
Em relação às Categorias Administrativas, tanto as Públicas ( 13,6 ) quanto as Privadas ( 13,1 ), apresentaram um desvio padrão inferior ao do Brasil ( 14,1 ).
Por sua vez, em relação às Organizações Acadêmicas, as Universidades ( 14,3 ) apresentaram um desvio padrão superior ao dos Centros Universitários ( 13,1 ), Faculdades ( 12,9 ) e Brasil como um todo, indicando uma dispersão um pouco maior nas notas das Universidades do que nas demais Organizações Acadêmicas.
Constata -se que existe diferença estatisticamente significativa ao nível de 95 \% nas médias das notas dos estudantes provenientes de Universidades em relação aos de Centros Universitários e Faculdades. 3.1.2 Estatísticas Básicas no Componente de Formação Geral A Tabela 3.3 apresenta as estatísticas básicas em relação ao componente da prova que avalia a Formação Geral dos estudantes concluintes.
Os alunos de todo Brasil obtiveram desempenho médio de 62,0.
Quanto à variabilidade, o desvio padrão das notas dos estudantes do Brasil, como um todo, foi 16,4.
A maior média foi obtida na região Sul ( 64,5 ), e a menor, na região Centro-Oeste ( 58,8 ).
As demais médias foram : 60,2 na região Norte e 62,0 nas regiões Nordeste e Sudeste.
Já o maior desvio padrão foi obtido na região Centro-Oeste ( 17,6 ), e o menor, na região Nordeste ( 15,4 ).
Os demais desvios padrões foram : 16,5 na região Norte ; 16,6 na região Sudeste, e 15,7 na região Sul.
A maior nota no Componente de Formação Geral da prova do ENADE/2014 foi 98,8, obtida por, pelo menos, um aluno na região Sul.
A menor nota máxima foi obtida na região Norte ( 97,6 ).
A mediana do Brasil, como um todo, foi 63,8, sendo a menor mediana encontrada na região Centro-Oeste ( 60,2 ) ; e a maior encontrada na região Sul ( 66,7 ).
A nota mínima nesta parte da prova foi zero nas regiões Sudeste e Centro-Oeste.
Nas demais regiões foi : 7,5 nas regiões Norte e Sul ; e 8,6 na região Nordeste.
Considerando -se as notas segundo Grande Região, observa -se que existe diferença estatisticamente significativa entre a maior média das notas do Componente de Formação Geral, obtida na região Sul ( 64,5 ), e a da região Centro-Oeste ( 58,8 ), a menor média.
O Gráfico 3.2 propicia a avaliação do desempenho dos estudantes do Componente de Formação Geral, a partir do histograma da distribuição das notas correspondentes.
A distribuição é unimodal, com moda em ( 60 ; 70 ], enquanto na prova, como um todo ( Gráfico 3.1 ), a moda foi alcançada dois intervalos abaixo ( 40 ; 50 ].
Destaca -se também o intervalo ( 70 ; 80 ] que possui valor muito próximo do modal.
Nota -se, ainda que, no Gráfico 3.2, as notas apresentam maior dispersão do que no Gráfico 3.1 ( distribuição das notas da prova ), confirmado pela comparação dos desvios padrões : 14,2 para a nota da prova como um todo e 16,4 para o Componente de Formação Geral.
Na Tabela 3.4 são apresentadas as informações referentes ao desempenho dos concluintes do Componente de Formação Geral, em diferentes agregações : Categoria Administrativa e Organização Acadêmica.
Observa -se que existe diferença estatisticamente significativa entre as médias por tipo de Categoria Administrativa.
A maior média foi obtida por estudantes de IES Públicas ( 67,0 ), com uma diferença estatisticamente significativa da obtida por estudantes de IES Privadas ( 59,8 ).
Considerando -se o tipo de Organização Acadêmica, nota -se que não há diferença estatisticamente significativa entre as médias dos alunos de Centros Universitários ( 59,6 ) e de Faculdades ( 59,2 ) ; mas a diferença é significativa entre a média obtida por esses e a dos estudantes de Universidades ( 63,9 ). 3.1.3 Estatísticas Básicas do Componente de Conhecimento Específico A Tabela 3.5 apresenta as estatísticas básicas referentes ao Componente de Conhecimento Específico da Área de Engenharia Ambiental.
A média do desempenho dos alunos do Brasil, como um todo, foi 42,2.
A maior média foi obtida na região Sul ( 48,4 ), e a menor, na região Centro-Oeste ( 38,9 ).
As demais médias foram : 39,4 na região Norte ; 40,7 na região Nordeste ; e 41,5 na região Sudeste.
Quanto à variabilidade das notas, o desvio padrão do Brasil, como um todo, foi 15,8, sendo o maior desvio padrão observado na região Sudeste ( 15,9 ), e o menor, na região Sul ( 14,7 ).
Os demais desvios foram : 15,2 na região Norte ; 15,1 na região Nordeste ; e 15,3 na região Centro-Oeste.
A mediana das notas dos estudantes de todo o Brasil foi 42,0.
A maior mediana ocorreu na região Sul ( 49,0 ), e a menor, na região Centro-Oeste ( 38,2 ).
As demais medianas foram : 40,0 nas regiões Norte e Nordeste ; e 41,0 na região Sudeste.
A nota máxima do Brasil, como um todo, foi 95,8, sendo obtida por, pelo menos, um aluno na região Sudeste.
As demais notas máximas foram : 80,0 na região Norte ; 87,2 nas regiões Nordeste e Sul ; e 84,7 na região Centro-Oeste.
A nota mínima foi diferente de zero apenas nas regiões Nordeste ( 5,0 ) e Sul ( 6,0 ).
Observa -se que existe diferença estatisticamente significativa da menor média das notas do Componente de Conhecimento Específico da região Centro-Oeste ( 38,9 ) em relação à região Sul ( 48,4 ), a maior média.
O intervalo de confiança para a média da região Norte ( 2,9 ) é maior do que de todas as outras.
Assim como os Gráficos 3.1 e 3.2, o Gráfico 3.3, apresentado a seguir, permite uma avaliação do desempenho de concluintes de Engenharia Ambiental em relação ao Componente de Conhecimento Específico com um histograma da distribuição das notas correspondentes.
Dentre as três distribuições apresentadas até aqui, esta é a mais concentrada nas notas baixas.
Esta também é uma distribuição unimodal e o grupo modal é o ( 40 ; 50 ], o mesmo grupo modal da prova como um todo e abaixo do grupo modal para a Formação Geral.
Destaca -se que o intervalo ( 30 ; 40 ] tem uma frequência muito próxima ao modal.
A Tabela 3.6 apresenta uma comparação dos resultados em relação à Categoria Administrativa e à Organização Acadêmica, agora levando em conta o desempenho de alunos do Componente de Conhecimento Específico da prova.
No que se refere à Organização Acadêmica, a maior média foi das Universidades ( 45,1 ), vindo a seguir a dos Centros Universitários ( 39,1 ) e, depois, a das Faculdades ( 37,9 ).
O maior desvio padrão, com valor superior ao do Brasil como um todo, foi o das Universidades ( 16,0 ).
As Universidades também obtiveram as maiores notas máxima ( 95,8 ) e mediana ( 45,3 ).
Por sua vez, as Faculdades obtiveram a maior nota mínima ( 1,5 ).
Já as Universidades e os Centros Universitários obtiveram nota mínima zero.
Quanto à Categoria Administrativa, observa -se um comportamento semelhante àquele da parte de Formação Geral e da prova como um todo, ou seja, existe diferença estatisticamente significativa entre as médias das IES Públicas ( 50,1 ) e IES Privadas ( 38,8 ).
Como nas análises anteriores, a maior média foi obtida por alunos de IES Públicas de ensino. 3.2 ANÁLISE DAS QUESTÕES OBJETIVAS Esta seção apresenta estatísticas para o conjunto de questões objetivas e uma análise gráfica de questões selecionadas dos componentes de Formação Geral ( 3.2.1 ) e de Conhecimento Específico ( 3.2.2 ).
São também apresentadas e comparadas as médias das subpopulações caracterizadas por Grande Região. 3.2.1 Componente de Formação Geral A Tabela 3.7 apresenta as estatísticas básicas relativas às oito questões objetivas do componente da prova que abrange a Formação Geral dos estudantes.
A média do Brasil foi 62,9.
A menor média foi encontrada na região Centro-Oeste ( 60,4 ), e a maior, na região Sul ( 65,0 ).
As demais médias foram : 60,7 na região Norte e 62,9 nas regiões Nordeste e Sudeste.
O desvio padrão do Brasil foi 20,1, sendo o maior desvio padrão encontrado na região Centro-Oeste ( 21,5 ), e o menor, na região Nordeste ( 19,1 ).
Os demais desvios foram : 19,9 na região Norte ; 20,3 na região Sudeste ; e 19,5 na região Sul.
As medianas ( 62,5 ), as notas máximas ( 100,0 ) e as notas mínimas ( 0,0 ) foram iguais para quase todas as regiões, excetuando da nota mínima na região Norte ( 12,5 ).
Observa -se que existem diferenças estatisticamente significativas ao nível de 95 \% no Componente de Conhecimento Específico entre as notas das Universidades em relação aos Centros Universitários e Faculdades. 3.2ANÁLISE DAS QUESTÕES OBJETIVAS Esta seção apresenta estatísticas para o conjunto de questões objetivas e uma análise gráfica de questões selecionadas dos componentes de Formação Geral ( 3.2.1 ) e de Conhecimento Específico ( 3.2.2 ).
São também apresentadas e comparadas as médias das subpopulações caracterizadas por Grande Região. 3.2.1 Componente de Formação Geral A Tabela 3.7 apresenta as estatísticas básicas relativas às oito questões objetivas do componente da prova que abrange a Formação Geral dos estudantes.
A média do Brasil foi 62,9.
A menor média foi encontrada na região Centro-Oeste ( 60,4 ), e a maior, na região Sul ( 65,0 ).
As demais médias foram : 60,7 na região Norte e 62,9 nas regiões Nordeste e Sudeste.
O desvio padrão do Brasil foi 20,1, sendo o maior desvio padrão encontrado na região Centro-Oeste ( 21,5 ), e o menor, na região Nordeste ( 19,1 ).
Os demais desvios foram : 19,9 na região Norte ; 20,3 na região Sudeste ; e 19,5 na região Sul.
As medianas ( 62,5 ), as notas máximas ( 100,0 ) e as notas mínimas ( 0,0 ) foram iguais para quase todas as regiões, excetuando da nota mínima na região Norte ( 12,5 ).
A Tabela 3.8 apresenta o índice de facilidade e o índice de discriminação ( ponto bisserial ) para \textbf{cada} uma das questões objetivas do Componente de Formação Geral.
Quanto ao índice de facilidade, foram usadas as seguintes cores para diferenciar o nível de dificuldade da questão :  Azul para as questões classificadas com índice muito fácil ( \&gt ; =0,86 ), verde para as questões classificadas com índice fácil ( 0,61 a 0,85 ), amarelo para as questões classificadas com médio ( 0,41 a 0,60 ), vermelho para as questões classificadas com difícil ( 0,16 a 0,40 ) e roxo para as questões classificadas com muito difícil ( \&lt ; =0,15 ).
Já quanto ao índice de discriminação, foram usadas as seguintes cores para qualificar a questão :  As questões classificadas com índice fraco receberam a cor vermelho ( \&lt ; =0,19 ), as classificadas com médio receberam a cor amarelo ( 0,20 a 0,29 ), as classificadas com bom receberam a cor verde ( 0,30 a 0,39 ) e as classificadas com muito bom ( \&gt ; =0,40 ) receberam a cor azul.
As questões objetivas do Componente de Formação Geral, segundo o índice de facilidade, foram assim avaliadas : das oito questões, nenhuma teve o índice de facilidade classificado como muito fácil.
Cinco questões foram tidas como fácil, por terem índice de acertos situado na faixa entre 0,61 e 0,85 ( de 61,0 \% a 85,0 \% de acertos ).
Duas questões foram classificadas com índice de facilidade médio, situando -se no intervalo entre 0,41 e 0,60, ou seja, houve entre 41,0 \% e 60,0 \% de acertos, enquanto uma questão foi classificada na categoria difícil, situando -se no intervalo entre 0,16 e 0,40.
Por fim, nenhuma das questões apresentou menos de 15 \% de acertos, razão pela qual seria classificada como muito difícil.
Como já comentado, para análise das questões objetivas, relativas à Formação Geral, segundo o poder de discriminação, utilizou -se o índice de discriminação ( ponto A Tabela 3.8 apresenta o índice de facilidade e o índice de discriminação ( ponto bisserial ) para \textbf{cada} uma das questões objetivas do Componente de Formação Geral.
Quanto ao índice de facilidade, foram usadas as seguintes cores para diferenciar o nível de dificuldade da questão :  Azul para as questões classificadas com índice muito fácil ( \&gt ; =0,86 ), verde para as questões classificadas com índice fácil ( 0,61 a 0,85 ), amarelo para as questões classificadas com médio ( 0,41 a 0,60 ), vermelho para as questões classificadas com difícil ( 0,16 a 0,40 ) e roxo para as questões classificadas com muito difícil ( \&lt ; =0,15 ).
Já quanto ao índice de discriminação, foram usadas as seguintes cores para qualificar a questão :  As questões classificadas com índice fraco receberam a cor vermelho ( \&lt ; =0,19 ), as classificadas com médio receberam a cor amarelo ( 0,20 a 0,29 ), as classificadas com bom receberam a cor verde ( 0,30 a 0,39 ) e as classificadas com muito bom ( \&gt ; =0,40 ) receberam a cor azul.
As questões objetivas do Componente de Formação Geral, segundo o índice de facilidade, foram assim avaliadas : das oito questões, nenhuma teve o índice de facilidade classificado como muito fácil.
Cinco questões foram tidas como fácil, por terem índice de acertos situado na faixa entre 0,61 e 0,85 ( de 61,0 \% a 85,0 \% de acertos ).
Duas questões foram classificadas com índice de facilidade médio, situando -se no intervalo entre 0,41 e 0,60, ou seja, houve entre 41,0 \% e 60,0 \% de acertos, enquanto uma questão foi classificada na categoria difícil, situando -se no intervalo entre 0,16 e 0,40.
Por fim, nenhuma das questões apresentou menos de 15 \% de acertos, razão pela qual seria classificada como muito difícil.
Como já comentado, para análise das questões objetivas, relativas à Formação Geral, segundo o poder de discriminação, utilizou -se o índice de discriminação ( ponto O Gráfico 3.4, para exemplificar, analisa o comportamento da questão de número 7 de Formação Geral.
Trata -se de uma questão considerada fácil, com relação a facilidade e a que obteve o maior índice de discriminação dessa parte da prova.
Neste gráfico, \textbf{cada} uma das cinco curvas representa o percentual de respostas em determinada alternativa da questão, em função do número de acertos dos estudantes nessa parte da prova ( Formação Geral/Múltipla Escolha ), antes de possíveis eliminações pelo critério do ponto bisserial.
A curva em vermelho bisserial ).
Nesta análise, as questões foram assim avaliadas : seis das oito questões apresentaram índice acima ou igual a 0,40 e, assim, foram classificadas com índice muito bom para esse grupo de estudantes.
As duas questões restantes tiveram índice de discriminação bom, com valor entre 0,30 e 0,39 e nenhuma questão teve nível médio ou fraco de discriminação para esse grupo de estudantes.
O índice de facilidade variou de 0,37 a 0,84, e o de discriminação, de 0,37 a 0,53.
As seis questões com índice de discriminação muito bom figuraram entre três níveis de dificuldade : quatro classificadas na categoria fácil ( questões 3, 4, 6 e 7 ) do índice de facilidade ; uma, na categoria médio ( questão 1 ) e uma, na categoria difícil ( questão 5 ).
Em particular, a questão 7 foi a que apresentou o maior poder discriminatório, com índice 0,53, porém foi considerada fácil em termos de facilidade, com uma proporção de 0,65 acertos.
O máximo de acertos foi alcançado pela questão 4 com um índice de facilidade de 0,84.
A questão de número 1 apresentou um índice de facilidade de 0,52, ou seja, um pouco mais da metade dos estudantes conseguiu resolvê -la, dentro do universo de participantes.
Seu índice de discriminação foi muito bom ( 0,42 ).
Já a questão 5 obteve índice de discriminação muito bom, 0,42, porém seu índice de facilidade foi difícil ( 0,37 ).
O Gráfico 3.4, para exemplificar, analisa o comportamento da questão de número 7 de Formação Geral.
Trata -se de uma questão considerada fácil, com relação a facilidade e a que obteve o maior índice de discriminação dessa parte da prova.
Neste gráfico, \textbf{cada} uma das cinco curvas representa o percentual de respostas em determinada alternativa da questão, em função do número de acertos dos estudantes nessa parte da prova ( Formação Geral/Múltipla Escolha ), antes de possíveis eliminações pelo critério do ponto bisserial.
A curva em vermelho corresponde à alternativa E, a correta para esta questão.
Assim, observa -se que entre os estudantes com menor número de acertos, nessa parte do exame, a situação mais frequente foi a escolha de uma das alternativas incorretas : a alternativa C ( em preto ) ou D ( em roxo ).
À medida em que o número de acertos aumenta, indicando desempenho melhor nessa parte da prova, aumenta concomitantemente a proporção de estudantes que selecionaram a alternativa correta E, atingindo 100 \% apenas para os estudantes com 8 acertos.
Essa análise permite verificar como a questão discriminou os grupos de desempenho, justificando o alto índice obtido na questão.
Os gráficos relativos às demais questões de Formação Geral constam do Anexo I. 3.2.2 Componente de Conhecimento Específico A Tabela 3.9 apresenta as estatísticas básicas em relação às questões objetivas do Componente de Conhecimento Específico da prova de Engenharia Ambiental por Grande Região.
A média do Brasil deste componente foi de 45,7.
A menor média foi observada na região Centro-Oeste ( 42,2 ), e a maior, na região Sul ( 51,8 ).
O desvio padrão de todo o Brasil foi 17,2, sendo o menor desvio padrão encontrado na região Sul ( 16,0 ), e o maior, na região Sudeste ( 17,3 ). ( 51,8 ).
O desvio padrão de todo o Brasil foi 17,2, sendo o menor desvio padrão encontrado na região Sul ( 16,0 ), e o maior, na região Sudeste ( 17,3 ).
A mediana de todo o Brasil foi 47,1, a mesma encontrada na região Sudeste.
A mediana foi menor ( 41,2 ) nas regiões Norte, Nordeste e Centro-Oeste ; enquanto foi maior na região Sul ( 52,9 ).
A nota máxima da prova foi 100,0, obtida nas questões objetivas do Componente de Conhecimento Específico, por, pelo menos, um aluno na região Sudeste.
Nas demais regiões, a nota máxima da prova foi : 88,2 nas regiões Norte e Centro-Oeste e 94,1 nas regiões Nordeste e Sul.
Em quase todas as regiões a nota mínima foi zero, exceto nas regiões Nordeste e Centro-Oeste ( 5,9 em ambas ).
A Tabela 3.10 apresenta os índices de facilidade e discriminação ( ponto bisserial ) das questões objetivas do Componente de Conhecimento Específico para os estudantes de Engenharia Ambiental.
Para facilitar a diferenciação das questões, usaremos as mesmas cores da Tabela 3.8 para as diferentes classificações dos índices de facilidade e de discriminação.
Dentre as questões objetivas da parte da prova relativa ao Componente de Conhecimento Específico, uma questão foi anulada pela Comissão, a de número 31.
Desse modo, a classificação, quanto ao índice de facilidade, foi estabelecida com base em 26 questões.
A partir dos índices obtidos, pode -se concluir que quase dois terços das questões objetivas da prova foi considerado, pelo menos, difícil : das 26 questões válidas, 16 foram classificadas como difícil ou como muito difícil ( a classificação modal ).
Duas questões foram classificadas como fácil e outras oito consideradas como médio.
Nenhuma questão foi classificada como muito fácil.
Já quanto aos índices de discriminação das questões objetivas do Componente de Conhecimento Específico da prova, tem -se como resultado a seguinte classificação : sete das 26 questões foram consideradas como boas, enquanto cinco delas tiveram índice de discriminação muito bom.
Assim, para quase a metade das questões – 12 em 26 – os índices de discriminação foram bom ou muito bom.
Dentre as demais, cinco delas foram classificadas como médio e outras nove como fraco, sendo 14, por conseguinte, a quantidade de questões nos dois patamares mais baixos de discriminação.
Constata -se, assim, que a prova – no que se refere ao Componente de Conhecimento Específico – possuía capacidade mediana de discriminar entre aqueles que dominam ou não o conteúdo.
Dentre as questões que alcançaram os maiores índices de discriminação, cinco delas, as de números 11, 17, 19, 23 e 32, foram classificadas com índice de discriminação muito bom, situando -se no intervalo de 0,40 a 0,42 do índice, sendo uma delas ( questão 23 ) classificada na categoria fácil, quanto ao índice de facilidade, e as outras quatro ( questões 11, 17, 19 e 32 ), na categoria médio.
A questão de número 26 foi a mais difícil dentre as questões específicas válidas, com baixo índice de facilidade, apenas 10,0 \% de acertos.
Essa questão apresentou poder discriminatório igualmente baixo, 0,13, o que comprova ter sido a mais difícil para os estudantes.
Destaca -se, também, a questão 35, com índice de facilidade 0,11, o que, em termos percentuais, corresponde a 11,0 \% de estudantes que responderam acertadamente.
Na questão 35, o índice de discriminação, 0,001, foi ainda mais baixo que o da questão 26.
Tais questões foram, portanto, pelo critério ponto bisserial, consideradas inadequadas.
Por isso, as questões 26 e 35 foram eliminadas do cômputo da nota final.
Além destas duas, as demais questões com índice de discriminação fraco, questões 10, 12, 22, 24, 27, 29 e 34 também não foram utilizadas no cômputo final das notas, num total de nove questões eliminadas.
A título de exemplo das análises do comportamento das questões objetivas, o Gráfico 3.5 analisa a questão 17 do Componente de Conhecimento Específico.
Esta questão foi classificada como médio.
Apresentou índice de facilidade 0,60, ou seja, 60,0 \% dos estudantes assinalaram acertadamente a opção B, correspondente ao gabarito.
Seu índice de discriminação foi igual a 0,42, classificado como muito bom, também sendo o maior índice discriminatório.
Neste gráfico, \textbf{cada} uma das cinco curvas representa o percentual de respostas em determinada alternativa da questão 17, em função do número de acertos dos estudantes nessa parte da prova, antes de possíveis eliminações de questões pelo critério do ponto bisserial.
A alternativa correta B, representada no gráfico pela curva em verde, foi escolhida em maiores proporções pelos alunos com desempenho melhor Já as alternativas incorretas, também denominadas distratores, foram selecionadas, principalmente, por aqueles com notas mais baixas.
Neste caso também a soma não é sempre 100 \% por causa das questões não respondidas ou com mais de uma opção marcada.
Aqueles com nota zero, na sua quase totalidade deixaram esta questão em branco ou marcaram mais de uma alternativa, comportamento considerado inválido.
A proporção de alunos que selecionou a resposta correta ( B ) aumenta gradativamente, chegando a atingir 100 \% para 18, 19 ou 21 acertos, enquanto a proporção dos que escolheram alternativas incorretas decai, a partir de quatro acertos, como função do número de acertos nessa parte da prova.
Os gráficos relativos às demais questões do Conhecimento Específico constam do Anexo I 3.3 ANÁLISE DAS QUESTÕES DISCURSIVAS Esta seção apresenta estatísticas selecionadas e histogramas das Questões Discursivas de Formação Geral ( 3.3.1 ) e Conhecimento Específico ( 3.3.2 ).
São também apresentadas e comparadas as médias de subpopulações caracterizadas por Grande Região. 3.3.1 Componente de Formação Geral As análises dos resultados de desempenho dos estudantes de Engenharia Ambiental nas duas questões discursivas relativas à Formação Geral, considerando -se as notas de conteúdo e de desempenho linguístico, encontram -se na Tabela 3.11 e no Gráfico 3.6.
Na Tabela 3.11, observa -se que a nota média nesse conjunto de questões foi próxima da obtida nas objetivas.
Os estudantes, de todo o Brasil, obtiveram, em Formação Geral, média 62,9 nas questões objetivas e 60,6 nas questões discursivas.
No entanto, pode -se notar um aumento do desvio padrão de 20,1 nas questões objetivas do Componente de Formação Geral dos alunos de todo o Brasil, para 22,1 nas questões discursivas do mesmo componente.
A maior média foi obtida na região Sul ( 63,6 ), e a menor, na região Centro-Oeste ( 56,4 ).
A mediana de todo o Brasil, neste componente, foi 65,0, a mesma nas regiões Nordeste e Sudeste ; na região Sul a mediana apresentou um valor maior, 68,0, e nas regiões Norte e Centro-Oeste apresentou um valor menor, 63,8 e 60,5, respectivamente.
A nota máxima ( 98,0 ) foi a mesma nas regiões Nordeste e Sudeste ; sendo 97,0 nas regiões Norte e Sul ; e 96,0 na região Centro-Oeste.
A nota mínima ( 0,0 ) foi a mesma em todas as regiões do Brasil.
O Gráfico 3.6 representa a distribuição das notas nas questões discursivas do Componente de Formação Geral.
A moda desta distribuição ocorre no intervalo ( 70 ; 80 ].
Destaca -se, também, o intervalo ( 60 ; 70 ], com distribuição próxima ao valor modal e acima de 20 \%.
Na sequência, os resultados verificados para \textbf{cada} uma das questões discursivas de Formação Geral serão apresentados, estabelecendo -se relações com os conteúdos abordados em \textbf{cada} uma delas.
Os comentários da Banca de docentes corretores, a respeito do observado na correção das respostas dos estudantes, suas impressões e conclusões serão apresentados junto à análise de \textbf{cada} questão.
Cumpre esclarecer que, tendo em vista que as questões discursivas de Formação Geral são padronizadas, ou seja, constam de todas as provas, os comentários da Banca são os mesmos para todas as carreiras acadêmicas, sendo direcionados a todos os estudantes que participaram do ENADE/2014.
A seguir, serão analisados os desempenhos da Área de Engenharia Ambiental nas duas questões discursivas de Formação Geral do ENADE/2014, comparando -se os resultados obtidos com comentários para \textbf{cada} questão. 3.3.1.1 Análise de Conteúdo da Questão Discursiva 1 do Componente de Formação Geral Os dados de Engenharia Ambiental, obtidos a partir das respostas à questão 1, encontram -se na Tabela 3.12 e no Gráfico 3.7.
Nessa questão – de melhor desempenho dentre as duas de Formação Geral – os alunos, de todo o Brasil, tiveram média 65,3.
A menor média para a questão 1 foi obtida na região Centro-Oeste ( 61,4 ), e a maior, na região Sul ( 69,8 ).
Quanto à variabilidade das notas, o desvio padrão de todo o Brasil foi 28,3.
O menor desvio padrão foi obtido na região Sul ( 26,4 ), e o maior desvio padrão foi obtido na região Norte ( 29,3 ).
A mediana do Brasil e da região Sudeste foram iguais, 70,0 ; as regiões Nordeste e Sul apresentaram mediana de valor maior ( 75,0 ) e as regiões Norte e Centro-Oeste apresentaram valor menor ( 65,0 ).
As notas máximas e mínimas da questão discursiva 1 foram as mesmas para todas as regiões do Brasil, respectivamente, 100,0 e 0,0.
O Gráfico 3.7 mostra a distribuição das notas na questão discursiva 1 do Componente de Formação Geral.
Observa -se que a maior frequência corresponde aos alunos que obtiveram nota dentro do intervalo ( 70 ; 80 ], correspondendo à moda da distribuição.
Destacam -se, também, os alunos que obtiveram nota no intervalo ( 90 ; 100 ], com a segunda maior frequência. 3.3.1.2 Comentários sobre a correção de Conteúdo das respostas à Questão Discursiva 1 O enunciado apresentado da Questão 1 era claro e foi redigido em linguagem acessível.
A questão abordava o tema da mobilidade urbana, relacionando -a aos ideais de desenvolvimento urbano sustentável estabelecidos pela ONU, do ponto de vista ecológico, cultural, político, institucional, social e econômico.
O enunciado apresentava uma tabela, sem referência à fonte de publicação, que detalhava percentualmente as modalidades de deslocamento da população urbana brasileira, com base em duas distinções : motorizado versus não motorizado ; coletivo versus individual.
O comando da questão solicitava, então, que o estudante desenvolvesse um texto dissertativo composto de duas partes : análise das consequências do transporte motorizado para o desenvolvimento sustentável e elaboração de duas propostas de intervenção pautadas pelo incentivo ao uso da bicicleta, meio de transporte não motorizado e, portanto, não poluente.
A temática dos efeitos da industrialização, com foco na produção de automóveis, e consequências para o meio ambiente é atual.
Além disso, a questão provoca a reflexão sobre a necessidade de adequação das cidades ao número de habitantes, no contexto da utilização dos meios de transportes de massa : metroviário, ferroviário, por ônibus.
Esperava -se, assim, que o concluinte recorresse a diferentes argumentos para analisar as consequências do transporte motorizado, explicitados pelo texto motivador, e propusesse duas ações de intervenção, as quais considerasse relevantes para o incremento ao uso de bicicleta, que, preferencialmente, deveriam manter coerência com a argumentação desenvolvida na primeira parte da resposta.
O conflito entre a busca pelo transporte individual e o coletivo é inerente a todas as grandes cidades do mundo.
No Brasil, esse conflito emerge com maior gravidade, como efeito do aumento da quantidade de veículos automotores ( automóveis e motocicletas ) e da baixa qualidade dos transportes coletivos.
Tal associação incentiva o uso de meios de transporte individual em detrimento da utilização dos coletivos.
O padrão de respostas era tecnicamente adequado ao enunciado e amplo o suficiente para abranger as variações de respostas apresentadas.
Deve ser ressaltado que a proteção ao ambiente ecologicamente equilibrado é tema conhecido dos concluintes dos cursos universitários, que deveriam possuir informações suficientes para desenvolver uma resposta abordando pelo menos dois dos aspectos listados no padrão de respostas como possibilidades de consequências ( item ‘ a ’ ) e duas propostas de ação de intervenção aceitáveis ( item ‘ b ’ ).
Segundo o padrão, concluintes poderiam desenvolver, quanto ao item ‘ a ’, os seguintes pontos : aumento da emissão de poluentes atmosféricos ; aumento da emissão de gases de efeito estufa ( CO2- dióxido de carbono, CO - monóxido de carbono, O3 - Ozônio ) ; aumento da poluição visual e sonora ; aumento da temperatura local e global ; aumento do consumo de combustíveis ; aumento de problemas de saúde ( cardíacos, respiratórios, dermatológicos ) ; aumento da frota de veículos, promovendo congestionamentos urbanos ; diminuição de áreas verdes ; desmatamento ; aumento das áreas de impermeabilização, resultando em enchentes, diminuição da infiltração da água e recarga de lençóis freáticos ; elevação dos custos de manutenção das cidades ( metroferrovias, rodovias, tratamento de água, limpeza da cidade, etc. ) ; necessidade de ampliação de vias trafegáveis ; necessidade de ampliação de áreas de estacionamento.
Em relação ao item ‘ b ’, os estudantes poderiam considerar os seguintes aspectos : construção de vias exclusivas para bicicletas ( ciclovias e ciclofaixas ) ; proposição de formas de integração entre o transporte por bicicletas, o metroviário e os ônibus coletivos, a fim de garantir segurança e conforto em momentos de adversidades climáticas e relevo \textbf{acidentado} ; pontos de aluguel e/ou empréstimos de bicicletas ; construção de bicicletários ; investimentos na segurança pública ; políticas de incentivo ao uso de bicicletas ( educação ambiental, qualidade de vida, saúde, propaganda ) ; implementação de políticas de crédito e de redução de custos das bicicletas.
Para o item ‘ a ’, foi encontrado um número expressivo de respostas adequadas.
As consequências mais citadas \textbf{estavam} relacionadas com a emissão de poluentes e o aumento dos engarrafamentos.
Houve críticas frequentes à qualidade do transporte público e à necessidade de locomoção rápida diante da distância das residências aos locais de trabalho.
Em tom de reivindicação, muitos estudantes defenderam solução para os problemas do transporte coletivo, quer metroviário, ferroviário ou por ônibus.
O uso da bicicleta também foi considerado por muitos como um meio de transporte benéfico à saúde, mas alguns apontaram a impossibilidade de sua utilização pelos idosos e deficientes.
No concernente ao item ‘ b ’, de modo quase uniforme citou -se a necessidade de ciclovias e ciclofaixas, bem como de se garantir segurança para os \textbf{ciclistas}.
A necessidade de manutenção das ciclovias e ciclofaixas existentes em algumas cidades foram também \textbf{lembradas}.
Com menos frequência também houve a indicação de maior conscientização da população, inclusive pela mídia, e pelo incentivo à aquisição desses veículos não motorizados, por meio da diminuição de impostos para redução do custo de aquisição.
Houve quem sugerisse, inclusive, uma política de doação de bicicletas.
Alguns mencionaram a necessidade de se assegurar integração das ciclovias com os demais modais, com a construção de bicicletários seguros diante das grandes distâncias a percorrer, notadamente nas metrópoles.
Também com base na dificuldade de o veículo não motorizado vencer as longas distâncias que o trabalhador deve percorrer nos centros urbanos para chegar ao local onde exerce a sua atividade produtiva, alguns alunos argumentaram contra a utilização da bicicleta como forma de resolução da poluição ambiental.
O principal equívoco cometido na sugestão de ações ( item ‘ b ’ ) foi o de mencionar a necessidade de melhoria do meio ambiente dissociado do uso de bicicletas, como se requeria no comando da questão.
Nesse sentido, foram encontradas respostas que indicavam a caminhada como opção.
Foram consideradas respostas fracas ( notas de zero a 30 ) aquelas meramente opinativas, ou que apenas transcreveram elementos do enunciado, sem qualquer acréscimo ou apreciação crítica.
As respostas medianas ( notas entre 35 e 70 ) foram as que não apresentavam algum desenvolvimento para os tópicos citados ( duas consequências e duas ações de intervenção ) ou por falta de alguns deles.
Já as respostas boas ( notas de 75 a 100 ) continham os tópicos considerados pelo padrão de respostas, com algum desenvolvimento, pelo menos, regular de argumentação.
As respostas dos estudantes revelaram um adequado conhecimento quanto aos efeitos do transporte motorizado para o ambiente e que a temática da ecologia parece ser tema bem compreendido, aparecendo nas respostas como uma preocupação para as futuras gerações.
Nesse sentido, algumas políticas públicas foram mencionadas apropriadamente para assegurar o desenvolvimento sustentável pelos respondentes, tais como : educação, habitação, saneamento, saúde, transporte e segurança.
No entanto, de modo geral, as respostas demostraram que os estudantes têm muita dificuldade na expressão escrita do pensamento, como se pode constatar pela avaliação do desempenho linguístico que ficou a cargo de uma banca específica, formada por profissionais da Área de Língua Portuguesa. 3.3.1.3 Análise de Conteúdo da Questão Discursiva 2 do Componente de Formação Geral A Tabela 3.13 mostra que o desempenho médio dos estudantes na questão discursiva 2 ( média 56,4 ) foi inferior ao obtido na questão discursiva 1 ( média 65,3 ).
A região Sul foi aquela cuja média, nessa questão, foi maior ( 57,8 ), e a de menor média foi a região Centro-Oeste ( 51,6 ).
Quanto à variabilidade das notas, o desvio padrão de todo o Brasil foi 30,4, superior ao obtido na questão discursiva 1 ( 28,3 ).
O maior desvio nessa questão foi obtido na região Norte ( 31,9 ), enquanto o menor foi obtido na região Sul ( 28,9 ).
A mediana de todo o Brasil foi 60,0, a mesma de quase todas as outras regiões, a exceção da região Sul ( 65,0 ).
As notas máximas ( 100,0 ) foram as mesmas em todas as regiões do Brasil, assim como as notas mínimas ( 0,0 ), sem exceção.
O Gráfico 3.8 mostra a distribuição das notas na questão discursiva 2 do Componente de Formação Geral.
Destaca -se o intervalo ( 70 ; 80 ], o modal, com uma frequência superior a 17 \%.
Os intervalos ( 40 ; 50 ], ( 60 ; 70 ] e ( 90 ; 100 ] são os intervalos com as maiores frequências, abaixo do modal e acima de 12 \%.
Nota -se, ainda, que as notas ficaram mais distribuídas ao longo dos intervalos em comparação à questão discursiva de número 1. 3.3.1.4 Comentários sobre a correção de Conteúdo das respostas à Questão Discursiva 2 A questão 2 abordava o tema da violência urbana a partir de um texto motivador recolhido em um site de um jornal brasileiro.
O fragmento relatava a prisão de três jovens por agressão a um outro jovem, menor de idade, guardador de carros.
A notícia trazia informações sobre o estado do jovem agredido e sobre as possíveis penas a serem aplicadas aos agressores.
O enunciado solicitava, então, que o estudante desenvolvesse um texto dissertativo composto de duas partes : análise de duas causas do tipo de violência descrito no texto e apresentação de dois fatores que contribuiriam para evitar a agressão relatada.
Como o texto motivador reduzia -se ao relato de um episódio de violência urbana, o enunciado não oferecia argumentos que ajudassem o estudante a elaborar uma análise das possíveis causas do fato.
As propostas a serem apresentadas, no sentido de se evitar em situações como a relatada, deveriam manter coerência com a análise das causas, desenvolvida na primeira parte da resposta.
O enunciado era claro e a matéria permitia identificar tratar -se de um conflito tipicamente urbano no qual se envolveram jovens de baixa renda em disputa por um território onde exerciam trabalho informal.
Sendo essa uma das hipóteses albergada sobre o amplo tema da violência urbana no qual se desenvolve o relato.
A temática da questão é bastante conhecida, não somente por debates no âmbito universitário, especialmente na Área de Ciências Sociais, bem como pela ênfase que é dada pela mídia onde assume contornos, muitas vezes, maiores do que o problema em si.
Assim, por tratar -se de tema bastante presente e debatido no cotidiano das pessoas, a questão pode ser considerada de baixa complexidade.
O padrão de resposta apresentava uma gama de possibilidades de causas da violência aceitas como corretas para o item ‘ a ’, tais como : problemas relacionados à educação ( baixa escolaridade, evasão escolar, qualidade da educação, distanciamento entre escola e a realidade social, tempo de permanência na escola ) ; desigualdades socioculturais ( gênero, etnia, economia, etc. ) ; desemprego e falta de qualificação profissional ; precariedade da segurança pública ; uso de drogas ; \textbf{desvalorização} da vida humana ; \textbf{banalização} da violência ; sensação de impunidade ; ausência de políticas sociais ; degradação da vida urbana ; desconhecimento ou desrespeito aos direitos humanos e constitucionais ; desestruturação familiar ; \textbf{desvalorização} de princípios éticos e morais.
No concernente ao item ‘ b ’, o padrão também previa a possibilidade de apresentação de diversos fatores que podem contribuir para evitar a violência : políticas de segurança mais efetivas ; políticas públicas de melhoria das condições socioeconômicas ; maior consciência cidadã e respeito à vida ; melhor distribuição de renda ; melhoria da educação ( aumento da escolaridade, redução da evasão escolar, qualidade da educação, aproximação entre a escola e a realidade social, aumento do tempo de permanência na escola ) ; aumento da oferta de emprego e melhoria da qualificação profissional ; medidas preventivas ao uso de drogas ; maior eficácia do sistema judiciário ; revisão da legislação penal ; valorização de princípios éticos, morais e familiares.
O enunciado da questão permitia que os estudantes, em suas respostas, pautassem o problema por diversos ângulos.
Dessa forma, o padrão buscou indicar diversas possibilidades a serem identificadas pelo corretor na leitura da resposta do estudante.
As notas foram atribuídas considerando o desenvolvimento da resposta e a indicação de ao menos dois dos itens exigidos pelo comando da questão nos itens ‘ a ’ ( duas causas ) e ‘ b ’ ( dois fatores ), dentre os considerados pelo padrão de respostas.
No entanto, o padrão exigia que as respostas \textbf{estivessem} afinadas com o perfil profissional previsto no Art. 3º da Portaria nº 255, de 02 de junho de 2014, onde se lê : Art. 3º No componente de Formação Geral serão considerados os seguintes elementos integrantes do perfil profissional : atitude ética ; comprometimento social ; compreensão de temas que transcendam ao ambiente próprio de sua formação, relevantes para a realidade social ; espírito científico, humanístico e reflexivo ; capacidade de análise crítica e integradora da realidade ; e aptidão para socializar conhecimentos em vários contextos e públicos diferenciados.
Nesse sentido, teses repressivas não conformadas com o texto constitucional e amparados pelos direitos humanos, não foram consideradas.
Questões como a majoração das penas, instituição de penas mais duras e \textbf{criminalização} da conduta dos menores desviados, só foram consideradas no caso de serem tratadas como menção a um debate que \textbf{está} presente, nunca como medida a ser aplicada de forma arbitrária e inconstitucional.
Também não foram previstas visões mais autoritárias e repressoras como a defesa de proibição do trabalho informal, como é o caso dos “ flanelinhas ”, e a defesa de não fornecer dinheiro aos “ guardadores ” de carro, identificado como esmolas por alguns.
Não foram consideradas respostas que meramente reproduziam o enunciado sem qualquer reflexão ou análise.
Aquelas que utilizaram os elementos do enunciado, mesmo que apenas transcritos, mas desenvolveram algum raciocínio pessoal ou indicaram causas e/ou fatores foram valoradas.
As causas apontadas com maior frequência para a violência foram : a ausência de educação e o desemprego, com soluções correlatas de educação profissional e de qualidade e abertura de postos de trabalho.
A questão das drogas também foi bastante mencionada, e a solução correlata mais indicada foi o acolhimento de usuários de drogas.
Outro tema bastante citado foi a dissolução das famílias e dos valores familiares.
Tendo em vista que o texto falava de moradores de rua, a ausência de moradia apareceu em muitas respostas, bem como a disputa territorial e a necessidade de sobrevivência.
Nesse sentido, as propostas apontavam para a necessidade de ajuda da Área de Serviço Social das prefeituras e atendimento psicológico gratuito.
Numa ótica mais repressora, surgiram críticas quanto à leniência das leis, acompanhadas da indicação da necessidade de atuação mais rigorosa do Poder Judiciário, como também da aplicação de penas maiores, negativa de esmolas, proibição do trabalho informal ( guardador de carros ), dentre outras.
Algumas respostas indicaram como causa o modelo capitalista, o consumismo e a ganância.
Outras indicaram o trabalho infantil e a correlata necessidade de aplicação do ECA ( Estatuto da Criança e do Adolescente ).
Em relação ao tema violência, inúmeras respostas apresentaram visões de intolerância quanto aos menores infratores, postulando a diminuição da idade para efeitos de \textbf{criminalização}, bem como defendendo a majoração de penas.
Essa perspectiva revela uma cisão social, de \textbf{certa} forma estimulada pelo tipo de divulgação da mídia quanto aos temas em foco, o que, sem dúvida, ficou evidenciado em algumas respostas.
A existência de inúmeras respostas pugnando pela maior repressão estatal e pelo aumento das penas, acrescido da busca pela diminuição da idade penal, revela que muitos estudantes não foram atingidos por uma formação superior que se exige humanista, ética e comprometida socialmente, que busque os fundamentos para resolver o problema da violência por meio de práticas democráticas e de inclusão social.
As respostas desse tipo também revelaram um alto grau de intolerância para com os moradores de rua, associando -os muitas vezes ao tráfico de drogas, à exploração indevida do espaço público, bem como a crimes como extorsão dos motoristas quando pedem dinheiro para guardar os veículos.
Nessa trilha, verifica -se lacuna na discussão de temas importantes para a formação de profissionais de nível superior tais como : sociodiversidade, multiculturalismo e violência ; tolerância/intolerância ; inclusão/exclusão.
Para aqueles que responderam em conformidade com o padrão de resposta verificou -se uma plena formação social, conclamando pelo auxílio não só do Estado, mas também da Sociedade no amparo daqueles que não têm as mínimas condições de sobrevivência e que \textbf{estão} nas ruas. 3.3.1.5 Análise de Língua Portuguesa das Questões Discursivas do Componente de Formação Geral Os dados de Engenharia Ambiental, obtidos a partir das respostas às questões discursivas do Componente de Formação Geral, no que tange à Língua Portuguesa, encontram -se na Tabela 3.14 e no Gráfico 3.9.
Nesse aspecto, os alunos, de todo o Brasil, tiveram média 59,8.
A maior média com respeito à Língua Portuguesa foi obtida na região Sul ( 63,0 ), e a menor, na região Centro-Oeste ( 55,8 ).
Quanto à variabilidade das notas, o desvio padrão de todo o Brasil foi 19,0.
O menor desvio padrão foi obtido na região Sul ( 17,6 ) e o maior desvio padrão foi obtido na região Centro-Oeste ( 19,9 ).
A mediana das notas de Língua Portuguesa foi 65,0 para o Brasil e para as regiões do Nordeste e Sudeste ; menor para a região Norte ( 62,5 ) e Centro-Oeste ( 60,0 ) ; e maior para a região Sul ( 67,5 ).
A nota máxima para todo o Brasil foi de 92,5, com, pelo menos, um aluno tirando essa nota na região Sudeste ; nas demais regiões a nota máxima foi : 90,0 nas regiões Norte, Nordeste e Sul ; e 87,5 na região Centro-Oeste.
Além disso, a nota mínima foi zero em todas as regiões do país.
O Gráfico 3.9 mostra a distribuição das notas de Língua Portuguesa do Componente de Formação Geral.
Observa -se que a maior frequência ( mais de 30 \% ) corresponde aos alunos que obtiveram nota no intervalo ( 60 ; 70 ]. 3.3.1.6 Comentários sobre a correção das respostas de Formação Geral com respeito à Língua Portuguesa Ao encaminhar as questões 1 e 2 na direção da produção de um texto dissertativo, esperava -se que o estudante utilizasse seus conhecimentos sobre o assunto e estruturasse seus textos de acordo com as exigências do registro formal próprio dessa situação comunicativa.
Essa configuração determina exigências quanto : à adequação da seleção vocabular, ao desenvolvimento do conteúdo, à estruturação sintática dos períodos, à organização lógica das ideias, à utilização de procedimentos de encadeamento textual e referenciação, à obediência às exigências morfossintáticas próprias da modalidade escrita da norma-padrão, ao respeito às regras ortográficas e às regras de acentuação gráfica.
O padrão de respostas utilizado na avaliação considerou os aspectos relevantes ao bom desempenho linguístico como competências distintas, de modo a permitir um mapeamento detalhado do domínio dos recursos disponíveis na Língua Portuguesa para a comunicação escrita formal.
Com base nesse objetivo, foram avaliados os seguintes aspectos : a ) Estruturação textual condizente com o gênero solicitado e o modo de organização textual expositivo adequado ao gênero – essa competência envolve : a estruturação sintática condizente com o padrão da modalidade escrita formal da língua portuguesa, de modo a garantir a clareza necessária ; a distribuição do conteúdo do texto em parágrafos, de modo a garantir a sua organização temática ; a utilização de operadores discursivos que contribuam para a progressão temática do texto, estabelecendo relações lógicas entre as ideias apresentadas, tanto do ponto de vista intrafrasal, como do interfrasal ; a utilização de procedimentos de referenciação lexical e pronominal que permitam a retomada de referentes textuais ; o respeito às regras de pontuação como fator de estruturação do período.
Espera -se, portanto, que o estudante recorra a procedimentos linguístico-discursivos para organizar seu texto, permitindo o encadeamento lógico entre suas partes de forma a garantir a progressão e a coerência textuais.
Isso significa que os seguintes procedimentos foram penalizados, de acordo com o padrão de respostas proposto : • a estrutura lógico-gramatical do texto fica comprometida pela elaboração de frases fragmentadas ; • sequência justaposta de ideias sem encaixamentos sintáticos, reproduzindo hábitos da oralidade ; • elaboração de frase com apenas oração subordinada, sem oração principal ; • emprego equivocado do conector ( preposição, conjunção, pronome relativo, alguns advérbios e locuções adverbiais ) que não estabeleça relação lógica entre dois trechos do texto e prejudique a compreensão da mensagem ; • emprego do pronome relativo sem a preposição, quando obrigatória ; • repetição ou substituição inadequada de palavras sem utilização dos recursos oferecidos pela língua ( pronome, advérbio, artigo, sinônimo ) ; • emprego inadequado dos pronomes relativos “ cujo(a) ” e “ onde ” ; • utilização inadequada dos sinais de pontuação que comprometa a clareza textual. b ) Respeito às convenções ortográficas da norma-padrão da Língua Portuguesa – essa competência envolve o domínio das regras de acentuação gráfica e da grafia padrão das palavras ( com ausência de abreviaturas próprias da linguagem da internet ), de acordo com as convenções estabelecidas pela legislação em vigor e consubstanciadas no Vocabulário Ortográfico da Língua Portuguesa, editado pela Academia Brasileira de Letras ( com aceitação da legislação anterior, no caso das regras relativas ao uso do hífen e da acentuação gráfica ).
Espera -se que o participante :  grafe corretamente as palavras ;  respeite as regras de acentuação gráfica ;  empregue maiúsculas em início de frase, em nomes próprios de pessoas, lugares ou instituições ;  evite abreviações como p/, vc, tb, pq, \textbf{tá}, né, usadas muitas vezes em escrita informal e na internet ;  obedeça às regras de separação de sílabas no final da linha. c ) Domínio dos diferentes aspectos morfossintáticos próprios da modalidade escrita formal da norma-padrão da Língua Portuguesa – essa competência envolve : a concordância nominal, a concordância verbal, a regência nominal, a regência verbal, a flexão nominal, a flexão verbal, a correlação entre os tempos verbais, a colocação pronominal e a utilização de sinais de pontuação que contribuam para a organização lógica da frase e do texto.
Espera -se que o participante :  flexione o verbo para estabelecer concordância de número com o sujeito da frase ;  flexione o artigo, o adjetivo e o pronome para concordar em número e em gênero com o substantivo a que se referem ;  observe a regência nominal e a verbal, utilizando a preposição adequada depois de um substantivo, um verbo ou um adjetivo ;  empregue adequadamente o acento grave indicador de crase nos casos em que se fizer necessário ;  obedeça às regras de colocação pronominal ( próclise e ênclise ), distintas dos hábitos da oralidade ou da escrita informal ; d ) Seleção vocabular adequada à modalidade escrita formal da Língua Portuguesa, exigida pela situação comunicativa – essa competência envolve a precisão na utilização do vocabulário relacionado à temática solicitada pela questão ; a ausência de marcas da oralidade, como termos de sentido de muito genérico ( “ coisa ”, “ negócio ”, “ você ” ) e termos de registros mais informais ( como gírias, jargões, frases feitas, ditados populares, termos regionais ).
Assim, espera -se que o participante respeite a adequação vocabular não utilizando gírias ou expressões coloquiais, evite repetição desnecessária de palavras e utilize um vocabulário mais formal, adequado ao texto de caráter dissertativo.
A escolha dessas competências para subsidiar o processo de avaliação apoia -se na concepção de que, no desempenho dos graduandos, a modalidade escrita tem apresentado uma intensa simplificação, originada no padrão da modalidade oral da Língua Portuguesa.
No caso do texto de base dissertativa, inscrito em um registro formal, a distância entre as duas modalidades é ainda maior, o que provoca situações de hipercorreção ( desvios provocados pela incorporação indevida de uma regra da norma-padrão ) e de truncamentos sintáticos ( estruturas frasais incompreensíveis devido à complexidade sintática própria da modalidade escrita ).
Observam -se, então, os seguintes aspectos que marcam essa distinção entre as duas modalidades, devido à excessiva simplificação da modalidade falada : a ) redução drástica de estruturas subordinadas, compensada pelo aumento na frequência de estruturas coordenadas e absolutas, por um lado, ou pela elaboração de estruturas truncadas pelo excesso de ideias sem a devida conexão subordinativa ; b ) redução no uso de conectores para expressar relações lógicas essenciais à construção do texto, substituídas pela exigência de inferência por parte do interlocutor para suprir a sua ausência ; c ) redução \textbf{cada} vez maior do uso do subjuntivo, ao lado da ampliação do uso do indicativo combinado a estruturas frasais coordenadas ou absolutas ; d ) empobrecimento do processo de referenciação, com a repetição exaustiva de pronomes ou nomes ; e ) simplificação extrema da marcação da categoria tempo na morfologia verbal ; f ) falta de domínio de vocabulário mais abstrato e de maior complexidade, essencial ao desenvolvimento do processo dissertativo ; g ) redução drástica no emprego da acentuação gráfica, processo intensificado pela divulgação imprecisa das mudanças promovidas pelo último acordo ortográfico.
Os aspectos macroestruturais da elaboração do texto não foram avaliados neste processo, para não penalizar duplamente os estudantes, já que a banca de formação geral, composta por profissionais de diferentes Áreas do conhecimento, encarregou -se da avaliação do conteúdo desenvolvido nas questões.
São eles : progressão temática, coerência na relação com os conhecimentos de mundo dos usuários da língua, inteligibilidade, atendimento ao solicitado no enunciado do ponto de vista do desenvolvimento do conteúdo, entre outros.
A grade de avaliação do desempenho linguístico considerou, portanto, três grandes grupos de competências, segundo os aspectos explicitados anteriormente : 1.
Aspectos ortográficos : domínio das convenções ortográficas : grafia de vogais e consoantes, uso de maiúsculas e minúsculas, emprego do hífen e acentuação gráfica ; 2.
Aspectos textuais : domínio dos procedimentos de estruturação textual do ponto de vista microestrutural : organização interna dos períodos, emprego de conectores para a articulação lógica entre os períodos e entre os parágrafos, emprego de marcas de referenciação lexical e pronominal ; utilização dos sinais de pontuação que contribuem para a organização lógica da frase. 3.
Aspectos morfossintáticos e vocabulares : domínio das regras de caráter morfossintático estabelecidas como modelares do ponto de vista da modalidade escrita formal da norma-padrão da Língua Portuguesa : concordância nominal e verbal, regência nominal e verbal, colocação pronominal, flexão nominal e verbal, correlação entre tempos e modos verbais, ausência de marcas de oralidade.
A seleção vocabular adequada à modalidade escrita formal da Língua Portuguesa foi incorporada a essa última competência, tendo em vista a intersecção entre as duas do ponto de vista das exigências do registro formal da modalidade escrita da norma-padrão.
Os resultados da avaliação correspondem aos seguintes aspectos observados em \textbf{cada} competência : Aspectos ortográficos : A correção foi realizada classificando os textos em cinco níveis, nível zero a quatro.
O desempenho dos estudantes ficou concentrado, majoritariamente, no nível 3 ( de 4 a 7 desvios ), devido, principalmente, aos desvios de acentuação.
O número de desvios de grafia foi reduzido.
Diferentemente do resultado do Enade 2013, aumentou a porcentagem de textos sem desvios ( enquadrados no nível 4 ) e diminuiu a porcentagem de textos enquadrados no nível 2.
Alguns textos, com pior desempenho, foram enquadrados no nível 1 ( de 8 a 12 desvios ), enquanto o nível zero foi atribuído a pouquíssimos casos, já que ele revela falta absoluta de domínio das convenções ortográficas.
Observou -se, portanto, que existe uma diferença muito grande de desempenho nos dois aspectos analisados : baixo índice de desvios ortográficos e grande índice de desvios de acentuação.
Em vários casos, ocorre ausência completa de acentuação gráfica.
Os resultados revelam, portanto, que a tendência dominante entre os universitários brasileiros é a eliminação da acentuação gráfica, provavelmente motivada pela vivência dos jovens relacionada aos aplicativos de comunicação via internet ( redes sociais e e-mails ).
Nesse tipo de comunicação, devido ao ritmo intenso de troca de mensagens, o uso de acentos gráficos foi praticamente abolido.
Outro fator que pode ter relação com essa tendência é a ausência de esclarecimento dos meios de comunicação, das autoridades e das escolas sobre as decisões do Acordo Ortográfico da Língua Portuguesa de 1990, gerando um estado de indefinição para os estudantes.
Os casos mais sistemáticos de eliminação do acento indicador da sílaba tônica são :  palavras proparoxítonas ( “ ridiculos ”, “ publicas ”, “ lideres ”, “ politicos ” ) ;  palavras paroxítonas terminadas em ditongo crescente ( “ homicidios ”, “ latrocinio ”, “ individuo ”, “ dependencia ” ) ;  palavras oxítonas ( “ ninguem ”, “ esta ”, “ ate ”, “ ai ” ).
Por outro lado, destaca -se o uso indevido do acento gráfico em determinadas palavras, como observado nas grafias * “ jornáis ”, * “ telejornaís ”, * “ propíciar ”, * “ medídas ”, * “ dígnidade ”, * “ cídades ”.
Quanto ao domínio das convenções relativas à grafia das palavras, observam -se desvios como : a hipercorreção pela escolha de “ e ” no lugar de “ i ”, por influência de hábitos da oralidade ( * “ enumeros ” por “ inúmeros ”, * “ entevenção ” por “ intervenção ” ) ; a eliminação do “ r ” marcador do infinitivo verbal ( * “ esta ” no lugar de “ \textbf{estar} ” ).
Outros casos de desvios de grafia relacionados à variação diastrática podem ser observados em *“estrupo ”, *“automovís ”, *“viensse ”, *“camihada ”.
São muito frequentes os seguintes desvios de caráter ortográfico, com repercussão morfossintática :  eliminação da marca de infinitivo ( -r- ) e substituição por acento agudo ( “ \textbf{está} ” no lugar de “ \textbf{estar} ” ) ;  confusão entre “ ão ” e “ am ” nas formas verbais ( “ invadão ” no lugar de “ invadam ” e “ estam ” no lugar de “ \textbf{estão} ” no presente do indicativo ; “ estaram ” no lugar de “ \textbf{estarão} ” no futuro do indicativo ) ;  confusão entre a grafia do verbo “ haver ” ( “ há ” ) e o artigo definido ou a preposição “ a ” ;  uso de hífen para separar pronome átono – tanto uso indevido quanto omissão ( no pretérito imperfeito e futuro do subjuntivo : “ evitar -mos ” no lugar de “ evitarmos ”, “ percebesse ” no lugar de “ percebe -se ” e vice-versa ) ; Observam -se, também, muitos casos de inadequação no uso da maiúscula : ausência de diferença entre a primeira letra e as outras, em início de período, principalmente dos estudantes que adotam a escrita em letra de imprensa ; utilização de maiúscula para destacar determinadas palavras-chave do texto, como “ Violência ”, “ Brasileiros ”, “ Fatores ”, “ Ozônio ”, “ Sustentável ”. para destacar determinadas palavras-chave do texto, como “ Violência ”, “ Brasileiros ”, “ Fatores ”, “ Ozônio ”, “ Sustentável ”.
Vale observar, também, que, ao contrário do que se esperava, não apareceram abreviaturas próprias do “ internetês ”, ou seja, dos hábitos de comunicação escrita adquiridos pelo uso de redes sociais e emails.
Aspectos textuais : Quanto a esses aspectos a correção também classificou os textos em cinco níveis ( zero a quatro ), em função da quantidade de erros apresentados.
O desempenho dos estudantes ficou concentrado, majoritariamente, nos níveis 3 e 2, devido à grande ocorrência de problemas de estruturação textual.
Foi muito baixo o número de textos que não apresentaram qualquer problema estrutural e, portanto, ficaram enquadrados no nível 4.
Alguns textos com pior desempenho foram enquadrados no nível 1, enquanto o nível zero foi atribuído a pouquíssimos casos, já que ele revela a existência de um texto sem articulação e com comprometimento do sentido.
Observou -se que a grande maioria dos estudantes não distribuiu as ideias em parágrafos, talvez devido ao pequeno número de linhas disponibilizadas para a resposta da questão ou, quem sabe, pela suposição de que não seria necessária essa divisão por não se tratar de um texto no modelo de uma redação dissertativo-argumentativa, como solicitado nos vestibulares.
Outro aspecto observado na estruturação textual foi a divisão em dois itens, provavelmente motivados pelo encaminhamento do enunciado das questões, que apresentavam dois pontos a serem detalhados : a questão 1 solicitava que os alunos dissertassem sobre as consequências do transporte motorizado ( a ) e apresentassem ações de intervenção por parte do poder público ( b ) ; a questão 2 solicitava que os alunos dissertassem sobre as causas da violência urbana ( a ) e os fatores para evitá -la ( b ).
Essa competência pode ser considerada como a mais problemática no que diz respeito ao desempenho linguístico dos estudantes, porque são muitos os problemas observados, desvios acumulados durante toda a formação escolar e que não se resolvem com um estudo autodidata, como acontece com regras ortográficas ou morfossintáticas : sequência justaposta de ideias sem encaixamentos sintáticos ; redução drástica de estruturas subordinadas, ao lado do aumento na frequência de estruturas coordenadas e absolutas ; redução no uso de conectores para expressar relações lógicas essenciais à construção do texto, substituídas pela exigência de inferência por parte do interlocutor para suprir a sua ausência ; emprego equivocado de operadores que não estabelecem relações lógicas coerentes entre ideias do texto ; emprego inadequado do pronome relativo ( com omissão da preposição ou a utilização de pronome inadequado, como “ onde ” ) ; repetição exaustiva de termos sem a utilização de procedimentos mais sofisticados de substituição ( hiperonímias, hiponímias, nominalizações, expressões metafóricas ) ; frases fragmentadas que comprometem a estrutura lógico-gramatical ; frases formadas apenas por oração subordinada, sem oração principal.
Um importante aspecto a destacar é o baixíssimo desempenho de uma parte dos estudantes em relação à estrutura formal do texto produzido, o que é extremamente preocupante ao se levar em conta que são graduandos em fase final de formação.
São frequentes os casos de desvios de estruturação frasal, com uso inadequado ou ausência de conectivos entre parágrafos e entre frases.
Em uma parte dos textos, falta um mínimo de textualidade e de domínio do registro padrão da língua.
Na verdade, observam -se relações linguísticas quase agramaticais, como as estabelecidas pela sequência de gerúndios sem o apoio de um ponto de partida para a organização das informações gramaticais e semânticas, ou seja, sem uma oração principal.
Quanto à utilização dos mecanismos de referenciação, deve -se destacar a ocorrência, em uma boa parte dos textos, de repetições de palavras ou expressões sem a utilização de termos sinônimos ou pronomes, como seria adequado.
Quanto à utilização dos sinais de pontuação, observou -se uma grande precariedade nos textos analisados.
É muito frequente a ocorrência de parágrafos sem marca interna de pontuação para separar os períodos.
Vale observar que não foi penalizada a ausência de vírgula para destacar locuções ou adjuntos adverbiais de pequena extensão deslocados de posição na frase, por ser um uso opcional.
São os seguintes os tipos de problemas encontrados : • vírgula : utilização de vírgula para separar o sujeito e o predicado ; ocorrência de apenas uma das vírgulas para separar uma palavra, uma expressão ou uma oração encaixada ; uso de vírgula no lugar do ponto para separar ideias que constituem períodos distintos ; ausência de vírgula para separar elementos de uma enumeração ; ausência de vírgula para separar oração adjetiva explicativa ou utilização inadequada para separar oração adjetiva restritiva ; • ponto e vírgula : utilização do ponto e vírgula no lugar de vírgula ; • ponto final : ausência de ponto final para separar períodos.
Aspectos morfossintáticos e vocabulares : Da mesma forma que nos aspectos anteriores os textos foram classificados em níveis de zero a quatro.
O desempenho dos estudantes ficou concentrado, majoritariamente, nos níveis 3 e 2, devido à grande ocorrência de problemas de regência e concordância.
O nível 4 foi atribuído a um número menor de textos.
Alguns textos com pior desempenho foram enquadrados no nível 1, enquanto o nível zero foi atribuído a pouquíssimos casos, já que ele revela a existência de um texto sem o respeito às mínimas exigências morfossintáticas da norma-padrão e com comprometimento do sentido.
Os resultados são muito transparentes em relação aos aspectos mais problemáticos do desempenho dos estudantes.
O desvio mais frequente, em relação à regência, é a falta do sinal indicativo da crase – isso revela que o usuário não tem consciência de que, sob a forma do termo “ a ”, existe a presença de uma preposição “ a ”, exigida pela regência do termo anterior.
Embora em outros exames, como o Enem, a falta de crase seja penalizada em acentuação, nesta correção esse desvio foi considerado no âmbito dos aspectos morfossintáticos.
Outro problema relacionado à regência verbal e à nominal, encontrado frequentemente nas questões, foi a ausência de preposição antes de pronome relativo, processo generalizado na modalidade oral da língua, em situações de registro informal.
Apesar da possibilidade de que essa alteração de regência se generalize no padrão escrito da Língua Portuguesa, como já \textbf{está} ocorrendo até em textos jornalísticos, o não emprego da preposição foi penalizado neste processo de avaliação.
Outro desvio muito frequente diz respeito aos processos de concordância verbal e de concordância nominal.
Quanto à concordância de número, observou -se ausência de marca ( com sujeito anteposto ou posposto ) ou uso indevido ( uso inadequado da marca de plural comandado pelo núcleo plural da locução adjetiva, apesar de o substantivo que funciona como núcleo do sintagma nominal \textbf{estar} no singular ).
Uma ocorrência que se destacou foi a ausência de acento circunflexo na forma plural do presente do indicativo dos verbos “ ter ” e “ vir ”, que foi considerada como um desvio na concordância verbal e não na acentuação gráfica.
Quanto à concordância de gênero, vários casos foram observados, normalmente no âmbito de sintagmas nominais longos, em que o adjetivo \textbf{está} afastado do substantivo.
Deve -se destacar uma ocorrência não observada no Enade/2013 : o aparecimento da marca de plural em verbos ou adjetivos relacionados a núcleos substantivos no singular, evidenciando um processo de hipercorreção.
Quanto à questão da colocação pronominal, foram poucos os casos observados.
Apesar de serem aspectos relacionados à oralidade, concluiu -se que, no registro escrito formal, a maioria dos estudantes já incorporou regras como a não introdução da frase por um pronome oblíquo e a próclise na presença de um termo atrator.
Não se adotou, entretanto, o padrão excessivamente formal descrito pelas gramáticas normativas em relação à posição do pronome oblíquo em locuções verbais, já que esse uso \textbf{está} muito distante da prática cotidiana, até em textos mais formais.
Quanto aos aspectos vocabulares, alguns tipos de inadequação foram observados : expressões da oralidade apareceram em algumas respostas, mas sem maior relevância do ponto de vista quantitativo ; seleção vocabular incompatível com o contexto, gerando situações de falta de inteligibilidade ; falta de domínio de vocabulário mais abstrato e de maior complexidade, essencial ao desenvolvimento do texto de base dissertativa.
Várias marcas de oralidade foram identificadas, embora não em alta frequência : o uso do pronome relativo “ onde ” como relativo universal, falta de artigo definido antes de substantivo, repetição de palavras por falta de vocabulário, reduções como “ \textbf{tá} ”, “ pra ”, “ pro ”, “ prum ”, expressões informais. 3.3.2 Componente de Conhecimento Específico Na parte da prova relativa às questões discursivas do Componente de Conhecimento Específico ( Tabela 3.15 ), observa -se que a média foi bem mais baixa do que para as questões discursivas do Componente de Formação Geral.
Enquanto no Componente de Formação Geral a média para estudantes de Engenharia Ambiental de todo o Brasil foi 60,6, na parte de Conhecimento Específico a média foi 22,5.
A maior média deste componente foi obtida pelos estudantes da região Sul ( 29,3 ), e a menor, pelos da região Nordeste ( 19,9 ).
Quanto à variabilidade das notas, o desvio padrão de todo o Brasil foi 15,7.
O maior desvio padrão foi encontrado na região Sul ( 16,2 ), e o menor, na região Centro-Oeste ( 14,5 ).
A maior nota máxima foi obtida na região Sul ( 86,7 ) e a menor na região Norte ( 61,7 ) ; as demais foram : 76,7 nas regiões Nordeste e Sudeste ; e 71,7 na região Centro-Oeste.
Além disso, a nota mínima ( 0,0 ) foi obtida por pelo menos um aluno em todas as regiões do Brasil.
A mediana do Brasil como um todo foi 21,7 ; nas regiões Norte, Sudeste e Centro-Oeste a mediana foi 20,0 ; na região Nordeste foi 18,3 ; e na região Sul foi 28,3.
O Gráfico 3.10 representa a distribuição das notas nas questões discursivas no Componente de Conhecimento Específico.
A moda desta distribuição ocorre no primeiro intervalo, [ 0 ; 10 ], intervalo que inclui os que deixaram a questão em branco, além dos que receberam zero para a resposta apresentada.
Nota -se que a distribuição é decrescente até o intervalo ( 80 ; 90 ].
A análise de \textbf{cada} uma destas questões será feita a seguir. 3.3.2.1 Análise da Questão Discursiva 3 do Componente de Conhecimento Específico Na questão 3, cujos resultados aferidos encontram -se descritos na Tabela 3.16, a média dos estudantes de todo o Brasil foi 22,5.
Nessa questão os estudantes obtiveram desempenho melhor que na questão discursiva 4 do Componente de Conhecimento Específico.
A menor média foi obtida pelos alunos da região Sudeste ( 20,2 ), enquanto a maior média foi obtida na região Sul ( 29,3 ).
Quanto à variabilidade das notas, o desvio padrão de todo o Brasil foi 20,7.
O maior desvio padrão foi obtido na região Norte ( 22,7 ), enquanto o menor foi obtido na região Sudeste ( 19,1 ).
A nota máxima, 100,0 pontos, foi alcançada por, pelo menos, um aluno das regiões Norte e Nordeste, enquanto nas demais regiões a nota máxima foi 95,0.
A mediana do Brasil como um todo, assim como na quase totalidade das regiões, foi 20,0, à exceção da região Sul, onde a mediana foi 30,0.
A nota mínima ( 0,0 ) foi obtida por, pelo menos, um estudante em todas as regiões do Brasil.
O Gráfico 3.11 mostra a distribuição das notas na questão discursiva 3, do Componente de Conhecimento Específico da Área de Engenharia Ambiental.
Sem levar em consideração que quase 17 \% dos estudantes deixaram a questão em branco, essa distribuição tem moda no intervalo ( 20 ; 30 ].
Destaca -se, também, o intervalo ( 10 ; 20 ] com distribuição muito próxima a moda. 3.3.2.2 Comentários sobre as respostas à Questão Discursiva 3 A questão possuía \textbf{certo} grau de subjetividade.
O enunciado fornecia um texto introdutório com uma explicação sobre resultados em relação aos reservatórios de água para consumo humano, do uso excessivo de agrotóxicos e de fertilizantes na agricultura.
Além disso, o pequeno texto introdutório mencionava o geoprocessamento como ferramenta útil para o acompanhamento do processo de eutrofização.
A questão visava avaliar a capacidade do estudante de correlacionar conhecimentos sobre geoprocessamento e sobre análises limnológicas.
Os conhecimentos exigidos para responder a essa questão – eutrofização – costumam ser apresentados na disciplina Ciência do Ambiente, de forma superficial, e, posteriormente, nas disciplinas específicas de Modelagem de Sistemas Ambientais, Tecnologias Ambientais e Poluição Ambiental, de forma mais aprofundada.
Os conhecimentos sobre geoprocessamento e imagens de satélite são estudados em Cartografia e em disciplina específica de Geoprocessamento.
Tais conteúdos \textbf{estão} previsto nas Diretrizes para Engenharia Ambiental, conforme Portaria Inep nº 243, de 02 de junho de 2014.
A questão exige, também, uma visão sistêmica do ambiente e um raciocínio crítico e criativo na identificação e resolução de problemas.
O enunciado da questão carecia de informações que delimitassem melhor o escopo da resposta esperada.
No item ‘ a ’, que solicitava que o estudante indicasse “ os parâmetros ”, não explicitava a quantidade mínima de parâmetros que deveriam ser citados e se era necessário que os estudantes mencionassem aspectos limnológicos e aspectos ligados ao geoprocessamento.
No item ‘ b ’, que solicitava explicação sobre como “ é possível constatar a ocorrência de eutrofização a partir de imagens de satélite ”, as respostas poderiam ser mais precisas se fosse informado que as imagens de satélite possuem uma resolução de 500 m a 1 km, já que a forma para constatar o fenômeno depende da resolução do satélite.
Além disso, na resposta ( explicação ) de muitos estudantes não havia nexo causal exclusivo com a eutrofização, excluindo outros fenômenos relacionados a outros tipos de poluentes.
A linguagem utilizada no enunciado era adequada ao jargão da Área de Engenharia Ambiental.
Apesar de ser considerada uma questão de média para difícil, o tempo e o espaço destinados para elaboração das respostas podem ser considerados suficientes para os que tivessem os conhecimentos necessários.
Em comparação com as demais questões discursivas de Conhecimento Específico o desempenho dos estudantes na questão pode ser considerado mediano.
No item ‘ a ’, o padrão de respostas esperada foi assim definido : “ O estudante deveria mencionar que, na estimativa do intervalo temporal a ser considerado na avaliação da eutrofização, poderia dispor de imagem de satélite e/ou de coletas e análises limnológicas ( citando, pelo menos, nitrogênio, fósforo e clorofila A ) desde quando o reservatório ainda não sofria com o carreamento superficial de nutrientes.
Esse intervalo depende da evolução temporal do uso e ocupação do solo na bacia hidrográfica estudada.
A estimativa de evolução da eutrofização fica prejudicada quando há carência de imagens de satélite e de dados de amostras de água ao longo do tempo. ” Nas respostas dadas a esse item, quase a totalidade dos estudantes citou apenas parâmetros do monitoramento limnológico convencional, com amostragens locais para determinação da qualidade da água.
Mesmo aqueles que na resposta ao item ‘ b ’ demonstraram bom conhecimento sobre o uso do sensoriamento para avaliação da eutrofização, no item ‘ a ’, citaram parâmetros como nitrogênio, fósforo e clorofila.
Um número menor de estudantes se preocupou em citar que esse monitoramento deve ser realizado desde o tempo em que o reservatório não \textbf{estava} poluído.
Ainda nesse item, poucos estudantes consideram o uso de imagens de satélite desde o tempo em que o reservatório não \textbf{estava} poluído como parâmetro para estimar o intervalo temporal.
Um número menor ainda citou tanto os principais parâmetros limnológicos ( nitrogênio, fósforo e clorofila A ) quanto o uso de imagens de satélite, além de considerar a temporalidade.
Foi comum, também, uma tentativa de parafrasear o texto do enunciado, citando apenas"acompanhamento dos agrotóxicos e fertilizantes", o que é uma resposta incorreta.
No item ‘ b ’, esperava -se que : “ O estudante deve mencionar que, de uma forma geral, a identificação da ocorrência da eutrofização e de sua evolução temporal é feita pela alteração da reflectância espectral do corpo hídrico na imagem de satélite.
É necessário que o estudante mencione que a análise deve ser iniciada anteriormente e prosseguir ao longo da ocorrência do fenômeno. ” Como resposta ao item ‘ b ’, alguns estudantes mencionaram como explicação de como é possível constatar a eutrofização por imagens de satélite apenas “ pela variação da cor ” ou “ pela variação da cor e do brilho do corpo d’água ”, sem falar em variação do verde ou do verde/azul.
Outros estudantes citaram o “ escurecimento do corpo d’água ou aumento de manchas causadas pelas algas ” e, ainda houve os que responderam que a constatação se dá pela : “ diminuição do espelho d’água ( pelo assoreamento inerente ao processo ) ” ; pelo “ acompanhamento da variação do uso do solo, verificando o aumento de áreas agrícolas ou urbanas ”.
Todas essas respostas \textbf{estão} apenas parcialmente corretas, pela impossibilidade de estabelecimento de nexo causal com o fenômeno da eutrofização.
A linguagem adotada pelos estudantes, em geral, não utilizou o jargão específico da Área de Engenharia Ambiental.
Apenas um número reduzido deles se expressou recorrendo a termos técnicos.
Além disso, foi bastante comum apresentarem uma redação com erros grosseiros de português.
De forma geral, o desempenho foi fraco ; com elevado percentual de questões em branco, de respostas anuladas e desconsideradas.
A grande quantidade de notas baixas pode ser justificada pela subdivisão do total de pontos da questão com maior valoração do item ‘ b ’ ( valia 70 do total de 100 pontos ), item que exigia o conhecimento da especificidade do uso do sensoriamento remoto para avaliação do processo de eutrofização, o que pode não ter sido abordado no curso.
Como já comentado, no item ‘ a ’, a quase totalidade dos estudantes citou apenas parâmetros do monitoramento limnológico convencional, com amostragens locais para determinação da qualidade da água.
Foi raro encontrar respostas completas que mencionassem os principais parâmetros limnológicos ( nitrogênio, fósforo e clorofila A ), o uso de imagens de satélite, além de considerar a questão da temporalidade.
No item ‘ b ’, a maioria das respostas foi imprecisa.
Um erro frequente foi afirmar que a constatação poderia ser feita observando o crescimento das algas ou das macrófitas.
Respostas como essas não atendem ao solicitado, já que a constatação deveria ocorrer “ a partir de imagens de satélite ” ( os satélites com maior resolução possuem 0,6 m de precisão, não permitindo diferenciar seres microscópicos ou pequenas macrófitas ).
Respostas afirmando que poderia ser observado o crescimento da espessura da camada de algas pela diminuição da penetração da luz, também não atendiam ao que foi solicitado.
Isto não é correto, pois os satélites só captam a luz refletida.
A diminuição da penetração da luz só pode ser avaliada com medições locais usando radiômetros ou disco de Secchi.
Apesar de o conhecimento exigido ser adequado ao previsto nas Diretrizes para Engenharia Ambiental, a especificidade do problema talvez tenha prejudicado o resultado, já que tema tão específico pode não fazer parte das ementas das disciplinas que abordariam o problema, em todos os cursos de Engenharia Ambiental.
É preciso que se reflita, também, sobre a necessidade de uma formação que prepare os futuros engenheiros ambientais para trabalharem com uma visão sistêmica do ambiente e um raciocínio crítico e criativo, na identificação e na resolução de problemas.
O baixo desempenho traz à tona a necessidade de inclusão nos currículos dos cursos de mais conteúdos específicos profissionalizantes, em disciplinas como : Modelagem de Sistemas Ambientais, Tecnologias Ambientais, Poluição Ambiental ; e, principalmente, de Cartografia e Geoprocessamento. 3.3.2.3 Análise da Questão Discursiva 4 do Componente de Conhecimento Específico A Tabela 3.17 contém as informações relativas à questão 4 do conjunto de questões do Componente de Conhecimento Específico.
O desempenho dos estudantes nessa questão foi inferior ao das questões 3 e 5.
A média geral do Brasil foi 21,0, sendo a menor média registrada na região Nordeste ( 17,8 ), e a maior, na região Sul ( 28,4 ).
A nota máxima ( 100,0 ) foi atingida por, pelo menos, um concluinte nas regiões Sudeste e Sul.
Nas regiões Norte e Centro-Oeste a nota máxima foi 80,0, e na região Nordeste foi 90,0.
A mediana em todo o Brasil, e em quase todas as regiões, foi 20,0, com exceção da região Sul, onde a mediana foi 30,0.
A nota mínima foi zero em todas as regiões.
O Gráfico 3.12 representa a distribuição de notas da questão discursiva 4, do Componente de Conhecimento Específico.
Essa distribuição é unimodal com moda no intervalo ( 10 ; 20 ].
A frequência de notas zero é alta, muito próxima da modal, assim como a categoria em branco, com aproximadamente 15 \%. 3.3.2.4 Comentários sobre as respostas à Questão Discursiva 4 Os conceitos exigidos pela questão, sobre aterro de resíduos e seu monitoramento, são abordados nos cursos de Engenharia Ambiental, de forma mais superficial, em alguma disciplina básica de Ciência do Ambiente e, de forma mais aprofundada, em disciplinas profissionalizantes como Resíduos Sólidos.
Disciplinas como Sistemas de Gestão Ambiental e de Avaliação de Impactos Ambientais abordam a temática do monitoramento ambiental e do plano de monitoramento associado ao controle de impactos ambientais.
Desta forma, a questão \textbf{está} adequada ao conteúdo previsto nas Diretrizes para Engenharia Ambiental ( Portaria Inep nº 243, de 02 de junho de 2014 ).
A necessidade de analisar uma situação e responder por meio de um texto dissertativo avaliava, também, a competência de análise com uma visão sistêmica do ambiente e o uso de raciocínio crítico e criativo para identificação e resolução de problemas.
A questão solicitava que o estudante dissertasse sobre “ importância da Elaboração do Plano de Monitoramento nas áreas de influência de um aterro de resíduos ”.
Para isso, ele deveria abordar aspectos enunciados em três itens orientadores que solicitavam : o conceito de Plano de Monitoramento, suas etapas nas várias fases do empreendimento e sua estrutura.
De modo geral, a questão era clara e redigida em linguagem adequada.
Dada a \textbf{certa} objetividade do que se esperava como resposta, pode -se considerar que o tempo e o espaço destinados para isso tenham sido adequados.
No entanto, a correção evidenciou que o enunciado deu margem a interpretações diferentes e, portanto, respostas diferentes da esperada pelo padrão definido para a correção.
Em comparação com as demais questões discursivas de Conhecimento Específico, a questão foi a que apresentou o pior desempenho dos estudantes, talvez devido a interpretações diferentes do que foi solicitado nos três itens do enunciado.
No item ‘ a ’, uma dificuldade detectada ocorreu por não \textbf{estar} explícito que o estudante deveria abordar o conceito de monitoramento com foco no que é “ exigido pela maioria das regulamentações ”, como delimitou o padrão de respostas.
Isso significa que a conceituação deveria \textbf{estar} ligada ao que normalmente é cobrado em um Termo de Referência para o desenvolvimento de um Estudo de Impacto Ambiental.
Nas respostas a esse item, boa parte dos estudantes abordou monitoramento ou controle para verificação de impactos ao meio ambiente ( degradação, comprometimento... ), causados pelo empreendimento.
Um número bem pequeno dos estudantes relacionou o monitoramento com processos previstos durante a elaboração do EIA do empreendimento.
No item ‘ b ’, muitos estudantes não compreenderam que deveriam abordar as etapas do empreendimento em que o monitoramento seria realizado, dentro da ideia de análise do ciclo de vida do empreendimento.
Segundo o padrão de resposta, o texto redigido pelos respondentes deveria mencionar as fases : diagnóstico, implantação, operação e encerramento.
Na maioria das respostas as etapas do monitoramento citadas foram o acompanhamento das águas subterrâneas e superficiais, do ar, do percolado, dos parâmetros geotécnicos.
Alguns associaram essas etapas às do licenciamento ambiental ( LP, PI e LO ).
Uma parcela menor citou as etapas associadas ao ciclo de vida ( diagnóstico e implantação, operação e encerramento ).
Em parte, essa variação ocorreu porque o comando não explicitou que se esperavam as etapas do empreendimento em que o monitoramento deveria ser realizado, dentro da ideia de análise do seu ciclo de vida.
Além disso, na bibliografia relacionada a resíduos sólidos e aterros sanitários a descrição do monitoramento do meio físico recebe mais atenção. 3.3.2.5 Análise da Questão Discursiva 5 do Componente de Conhecimento Específico A Tabela 3.18 contém as informações relativas à questão 5 do conjunto do Componente de Conhecimento Específico.
O desempenho dos estudantes de todo o Brasil nesta questão foi mais alto dentre as três questões discursivas desse componente.
A nota média dos estudantes de todo o Brasil foi 23,9.
A maior média foi registrada na região Sul ( 30,2 ), enquanto a menor média foi registrada na região No item ‘ c ’, que solicitava a estrutura do Plano de Monitoramento, o padrão de respostas delimitou que, no mínimo, a resposta deveria conter o que um plano desse tipo deve apresentar para garantir efetividade de controle, levando em conta a variação espaço-temporal.
Para isso, a resposta deveria “ considerar : tipo de amostra e os parâmetros a serem monitorados, a \textbf{localização} da amostragem, a periodicidade das coletas, a técnica de coleta, preservação e análises ”.
No entanto, muitos dos que responderam a esse item tiveram compreensão diversa do que seria apresentar a estrutura do plano de monitoramento.
A maioria dos estudantes não conseguiu descrever a estrutura do Plano de Monitoramento, sendo que alguns citaram que seria estruturado conforme o Termo de Referência do órgão de controle ambiental.
Uma parcela ainda menor citou um ou mais aspectos da estrutura mínima.
As respostas encontradas indicam que a solicitação do item não foi clara, por não explicitar que como estrutura do Plano de Monitoramento o estudante deveria citar o que esse deve conter para garantir a efetividade de controle, levando em conta a variação espaço-temporal.
A linguagem adotada pelos estudantes, em geral, não utilizou o jargão específico da Área de Engenharia Ambiental.
Apenas um número reduzido de estudantes se expressou com precisão.
Como já comentado na análise da questão 3, foram observados muitos erros grosseiros no uso da norma culta da língua portuguesa.
De forma geral, os estudantes apresentaram um desempenho fraco.
O percentual de respostas em branco e completamente erradas foi elevado.
Os resultados mostram que, a formação dos estudantes que participaram do exame de 2014 não garantiu aprendizagem significativa dos conteúdos básicos de Ciência do Ambiente e, também, daqueles abordados nas disciplinas profissionalizantes de Resíduos Sólidos, de Sistemas de Gestão Ambiental e de Avaliação de Impactos Ambientais, principalmente no que diz respeito a integração dos conhecimentos abordados nessas disciplinas. 3.3.2.5 Análise da Questão Discursiva 5 do Componente de Conhecimento Específico A Tabela 3.18 contém as informações relativas à questão 5 do conjunto do Componente de Conhecimento Específico.
O desempenho dos estudantes de todo o Brasil nesta questão foi mais alto dentre as três questões discursivas desse componente.
A nota média dos estudantes de todo o Brasil foi 23,9.
A maior média foi registrada na região Sul ( 30,2 ), enquanto a menor média foi registrada na região A nota mínima ( 0,0 ), a mediana ( 20,0 ) e a nota máxima ( 100,0 ) foram as mesmas em quase todas as regiões ; exceto na região Sul, onde a mediana foi 30,0, e na região Centro-Oeste, onde a nota máxima foi 95,0.
O Gráfico 3.13 representa a distribuição das notas da questão discursiva 5 do Componente de Conhecimento Específico.
Nota -se que a frequência de estudantes que deixaram a questão em branco é a modal e bastante próxima da frequência de notas no intervalo ( 10 ; 20 ].
Destaca -se que, a partir do intervalo ( 10 ; 20 ] a distribuição é decrescente até o intervalo ( 90 ; 100 ]. 3.3.2.6 Comentários sobre as respostas à Questão Discursiva 5 A questão, como a anterior ( questão 4 ), exigia conhecimentos relativos a Resíduos Sólidos.
Nessa, a problemática dos resíduos se articulava com conhecimentos sobre aproveitamento energético.
A questão solicitava que o estudante discorresse sobre “ O aproveitamento do metano para a produção de energia elétrica ”, orientado por três itens.
O primeiro, item ‘ a ’, solicitava que o estudante descrevesse o processo de formação do metano a partir de resíduos sólidos ; o item ‘ b ’ pedia um modelo para transformação desse gás em energia elétrica ; por fim, no item ‘ c ’, o estudante deveria apresentar uma possibilidade de comercialização dos gases gerados em aterros sanitários.
Com caráter apenas motivacional, o enunciado citava a lei que instituiu a Política Nacional de Resíduos Sólidos no texto introdutório, já que a resposta esperada não se articulava explicitamente com o texto.
O enunciado era claro, redigido em linguagem pertinente à Área.
Para o estudante de Engenharia Ambiental, os comandos dos três itens eram de fácil compreensão e demandavam respostas bastante objetivas.
Pela objetividade do que era esperado, o tempo e o espaço destinados à resolução da questão foram adequados.
Em comparação com as demais questões discursivas de Conhecimentos Específicos, nessa questão foi encontrado um número maior de notas boas, apesar de uma distribuição semelhante nas demais faixas.
Talvez isso se deva à maior clareza do enunciado e a comandos que demandavam respostas mais objetivas.
No item ‘ a ’, o estudante deveria mencionar que a matéria orgânica proveniente dos resíduos ( lixo doméstico, esgoto, fezes de animais em pastos ) sofre degradação em meio anaeróbio, quando ocorre fermentação.
Entre os compostos formados \textbf{estão}, basicamente, metano e dióxido de carbono.
Nas respostas a esse item, boa parte dos estudantes mencionou decomposição ( ou degradação, ou deterioração... ) dos resíduos e que esse processo gera gás carbônico ( ou dióxido de carbono ) e metano.
Alguns citaram que a degradação é anaeróbica, ou que a degradação se inicia de forma aeróbica até o recobrimento, mas que, depois, é anaeróbica.
Alguns, em menor quantidade, afirmaram que a degradação ocorre apenas com a fração orgânica dos resíduos, pela ação de micro-organismos ( bactérias ou arqueas ), por meio de um processo fermentativo ( ou através das etapas de hidrólise, acidogênese, acetogênese e metanogênese ).
Outros, de forma incorreta, afirmaram que o metano era oriundo de um processo aeróbico ou que era formado pelo chorume.
Na resposta ao item ‘ b ’ o estudante deveria citar que o modelo básico envolve sistema de captação e tratamento de gás, turbinas e/ou motores/geradores, e que esses equipamentos devem ser adequados ao volume de gás gerado.
Boa parte dos estudantes apresentou um modelo de transformação do metano em energia elétrica, não necessariamente com todos os detalhes, que incluía : a estrutura de captação, a drenagem e a canalização dos gases, um sistema de tratamento para a retirada da umidade do gás sulfídrico e de siloxanes, mas deixaram de mencionar o principal componente do modelo, que envolve um motor e um gerador, ou um motor que gera vapor e aciona um gerador.
Respostas desse tipo demonstraram conhecimento maior da área de Resíduos Sólidos e menor de Recursos Energéticos.
Já outra parte razoável dos estudantes só mencionou o principal componente do projeto ( turbinas e/ou motores geradores ), sem mencionar a estrutura de captação e tratamento.
Nos modelos apresentados foram encontrados os seguintes exemplos do conjunto motores geradores : caldeira/vapor/gerador, motor ciclo Otto, turbina, microturbina, termelétrica, térmica a vapor, dentre outros.
Desses, o mais adequado seria o motor ciclo Otto com gerador ou em microturbina, pela melhor eficiência.
Alguns mencionaram, de forma pouco precisa, “ máquina de produção de energia ” ou “ usina de produção de energia ”.
Outros escolheram biodigestores, o que é conceitualmente errado em modelos para transformação do metano em energia.
Houve, ainda, quem aconselhasse termelétricas que queimam o lixo e geram energia, o que também \textbf{está} errado, já que as térmicas a biomassa não são um modelo para gerar energia elétrica com o metano.
No item ‘ c ’, o estudante deveria mencionar que a comercialização pelo Mecanismo de Desenvolvimento Limpo ( MDL ) é realizada por meio da venda de crédito de carbono/redução certificada de emissões.
Esse item parece ter sido mais difícil para os estudantes.
Boa parte deles demonstrou não ter conhecimento do Mecanismo de Desenvolvimento Limpo ( MDL ) introduzido pelo Protocolo de Kyoto.
Como resposta citaram exemplos de venda para indivíduos ou de uso direto da energia gerada, o que não pode ser realizado via MDL.
Alguns indicaram que a comercialização seria realizada através de Reduções Certificadas de Emissões, ou simplesmente créditos de carbono.
Poucos registraram que uma tonelada de dióxido de carbono ( CO2 ) corresponde a um crédito de carbono.
Raro, também, foi mencionarem o potencial de aquecimento global e o uso da equivalência dos gases de efeito estufa em termos de tonelada de carbono ( a questão do carbono equivalente ).
Apenas, um número reduzido de estudantes se expressou com correção e \textbf{detalhamento}, usando termos técnicos de forma adequada.
Cabe registrar, mais uma vez, que os erros grosseiros no uso da Língua Portuguesa foram frequentes.
Nessa questão, uma parcela pequena das notas foram maiores ou iguais a 50, considerando o total de estudantes.
O percentual de respostas em branco foi maior do que nas questões 3 e 4.
Também foi mais elevado o percentual de respostas anuladas ou desconsideradas.
O melhor desempenho dos alunos ocorreu na resposta ao item ‘ a ’.
No item ‘ b ’ o desempenho foi ruim, o que mostra a necessidade de mais atenção dos cursos e/ou dos estudantes aos conteúdos profissionalizantes específicos de Tecnologia Ambiental e Recursos Energéticos.
Nas respostas dadas ao item ‘ c ’ observou -se que boa parte dos estudantes não tem conhecimento razoável sobre o Mecanismo de Desenvolvimento Limpo ( MDL ), introduzido pelo Protocolo de Kyoto, o que é uma deficiência grave para quem vai trabalhar na área de meio ambiente.
Não houve muita diversidade nas soluções que pudessem ser causadas por diferentes interpretações.
A variação das respostas ocorreu devido a diferentes graus de conhecimento sobre o assunto.
O baixo desempenho, além de evidenciar a necessidade de mais atenção aos conteúdos ( nas áreas de Resíduos Sólidos, de Tecnologia Ambiental e de Recursos Energéticos ), pode refletir falta de integração dos conhecimentos abordados em diferentes disciplinas, de modo a tornar o profissional capaz de mobilizá -los adequadamente na resolução de uma situação-problema.
Os frequentes erros de português mostram a formação deficiente da Educação Básica, que antecede o ensino superior, e que esse precisa enfrentar.

% matched lemmas: acidentado, banalização, cada, certo, ciclista, criminalização, desvalorização, detalhamento, estar, lembrar, localização
\end{textsample}
